% ============================================================================
% CHAPITRE 1 — Présentation de l'organisme & du sujet
% Structure alignée sur la table des matières cible
% ============================================================================
\openchapter{Présentation de l'organisme \& du sujet}

\section{Introduction}

Ce premier chapitre présente le cadre général dans lequel s'est déroulé ce stage de fin d'année. Il commence par une présentation de l'organisme d'accueil afin de situer l'environnement professionnel dans lequel le projet a été réalisé. Il expose ensuite le contexte et l'organisation du projet, avant de détailler le sujet de stage, sa problématique, la méthode de gestion adoptée, l'équipe impliquée et le planning de réalisation.

\section{Présentation de l'organisme d'accueil}

\textbf{Omni Networks} est une entreprise fondée à Agadir par une équipe d'ingénieurs spécialisée dans le développement de logiciels. Forte de plus de 20 ans d'expérience dans le domaine du développement de solutions métiers, l'entreprise met son expertise au service des entreprises afin de répondre à leurs différents besoins en matière de gestion et de fonctionnement.

L'activité de Omni Networks repose principalement sur la conception et le développement de solutions logicielles adaptées aux besoins de ses clients. L'entreprise accompagne ainsi ses clients dans l'évolution et l'amélioration de leurs outils et processus, quels que soient leur taille ou leur domaine d'activité.

Animée par une démarche orientée vers l'innovation et l'amélioration continue, Omni Networks s'appuie sur trois principes essentiels : \textbf{créer, innover et améliorer}. Cette approche lui permet de développer et de faire évoluer des produits de gestion destinés à accompagner durablement les entreprises dans leurs activités.

\section{Contexte et organisation du projet}

Le projet réalisé au cours de ce stage s'inscrit dans le développement d'une solution de gestion destinée à une association dont les activités reposent principalement sur les cotisations de ses abonnés et sur différentes sources de financement. L'association réalise notamment des projets à caractère social, tels que le forage de puits, ainsi que la construction et la rénovation de mosquées, de maisons pour veuves et d'orphelins. Les projets peuvent être financés par les cotisations des abonnés, les contributions des membres et des bienfaiteurs, ainsi que par des dons de l'État ou d'autres associations.

L'organisation de l'association repose sur plusieurs catégories de membres, parmi lesquelles le président, le vice-président, le trésorier, le vice-trésorier, le secrétaire général, le vice-secrétaire général et les conseillers. Les membres du bureau sont également considérés comme des abonnés de l'association.

Pour chaque projet, le bureau constitue un comité composé de membres de l'association. Ce comité est chargé d'assurer le suivi du projet de son lancement jusqu'à sa clôture. Selon leurs responsabilités et leurs droits, les membres du comité peuvent enregistrer les différentes opérations liées au projet, notamment les dons des adhérents, les paiements, les dépenses ainsi que les documents justificatifs tels que les reçus et les factures.

Dans ce contexte, le projet consiste à mettre en place une application permettant de centraliser ces informations et d'organiser leur gestion selon les rôles des différents utilisateurs. La solution doit également permettre de conserver les justificatifs associés aux opérations, de faciliter la préparation des rapports de clôture des projets et de permettre aux abonnés de consulter leur propre historique, comprenant leurs abonnements, leurs dons et les éventuelles sommes restant dues.

\section{Présentation du sujet de stage}

\subsection{Titre du projet}

\textbf{Conception et développement d'une plateforme web de gestion des abonnés et des projets d'une association}

\subsection{Problématique}

La gestion des abonnés, des cotisations et des projets d'une association implique le suivi d'un grand nombre d'informations : paiements des abonnements, dons, dépenses, factures, reçus et autres pièces justificatives. À cela s'ajoute la nécessité de suivre les projets individuellement et de permettre aux différents membres d'intervenir selon leurs responsabilités.

Le principal enjeu du projet est donc de centraliser ces informations au sein d'une même plateforme, tout en garantissant une gestion adaptée aux rôles des utilisateurs et une traçabilité des opérations réalisées. La solution doit également faciliter le suivi financier des projets, la conservation des documents justificatifs et la préparation des rapports de clôture.

Une autre problématique concerne les abonnés eux-mêmes. Chaque abonné doit pouvoir accéder à son historique afin de consulter ses abonnements, ses dons ainsi que les éventuelles sommes restant dues à l'association.

La problématique générale peut ainsi être formulée de la manière suivante :

\begin{quote}
\textit{Comment concevoir et développer une plateforme web permettant de centraliser la gestion des abonnés et des projets d'une association, d'assurer la traçabilité des opérations financières et de faciliter l'accès aux informations selon les rôles des différents utilisateurs ?}
\end{quote}

\subsection{Méthode de gestion de projet}

Compte tenu de la nature du projet et du fait que sa réalisation a été effectuée individuellement et à distance, une approche itérative et progressive a été adoptée. Le développement a été organisé en plusieurs étapes, allant de l'analyse des besoins et de la définition des fonctionnalités à la conception, au développement, aux tests et aux corrections.

Cette organisation a permis d'avancer progressivement sur les différentes fonctionnalités de l'application, tout en intégrant les ajustements nécessaires au fur et à mesure de l'avancement du projet. Les échanges réguliers avec le superviseur de stage ont permis de valider les orientations prises et de recueillir les remarques nécessaires à l'amélioration de la solution.

\subsection{Équipe de projet}

La réalisation de ce projet a été effectuée individuellement et à distance. En tant que stagiaire développeur, j'ai assuré l'ensemble des activités liées à la réalisation de la solution, notamment l'analyse des besoins, la conception, la modélisation, le développement, les tests ainsi que la correction des anomalies.

Le projet a néanmoins bénéficié d'un encadrement à deux niveaux. Le superviseur de stage a assuré le suivi professionnel du projet, en apportant ses orientations et ses recommandations tout au long de sa réalisation. Le tuteur académique a, quant à lui, assuré le suivi pédagogique du stage et veillé au bon déroulement du travail dans le cadre de la formation d'ingénieur.

L'organisation à distance a reposé principalement sur des échanges réguliers avec l'encadrement, permettant de suivre l'avancement du projet, de discuter des choix réalisés et de valider progressivement les différentes étapes.

\subsection{Planning du projet}

Le projet a été réalisé progressivement selon plusieurs phases, depuis l'analyse initiale des besoins jusqu'à la finalisation de l'application. Les principales étapes sont les suivantes :

\begin{enumerate}
    \item \textbf{Analyse des besoins} : étude du fonctionnement de l'association, identification des acteurs et définition des besoins fonctionnels et techniques.
    \item \textbf{Conception} : modélisation du système, définition de l'architecture de la solution et conception de la base de données.
    \item \textbf{Développement} : implémentation progressive des différentes fonctionnalités de l'application et intégration des composants du système.
    \item \textbf{Tests fonctionnels et techniques} : réalisation de tests sur les fonctionnalités développées afin de vérifier leur bon fonctionnement dans différents scénarios d'utilisation et d'identifier les éventuelles anomalies.
    \item \textbf{Correction et amélioration} : correction des anomalies détectées lors des tests et amélioration progressive des fonctionnalités et de l'interface.
    \item \textbf{Finalisation} : stabilisation de l'application, amélioration de l'expérience utilisateur et préparation de la documentation du projet.
\end{enumerate}

Ces phases n'ont pas été strictement séquentielles : conformément à la démarche itérative décrite plus haut, le développement, les tests et les corrections se sont recouverts d'une itération à l'autre, chaque nouveau module fonctionnel étant testé et corrigé avant que le suivant ne soit abordé. Le récapitulatif suivant associe à chaque phase son objectif principal et le livrable correspondant.

\begin{center}
\small
\begin{tabular}{@{}p{3.6cm}p{4.6cm}p{5.2cm}@{}}
\toprule
\textbf{Phase} & \textbf{Objectif principal} & \textbf{Livrable} \\
\midrule
Analyse des besoins & Comprendre le fonctionnement de l'association & Liste des acteurs et des besoins fonctionnels \\
\addlinespace[2pt]
Conception & Structurer la solution avant le codage & Modèle de données, architecture technique \\
\addlinespace[2pt]
Développement & Implémenter les modules fonctionnels & API REST et interfaces web opérationnelles \\
\addlinespace[2pt]
Tests fonctionnels et techniques & Vérifier le comportement attendu & Anomalies identifiées et documentées \\
\addlinespace[2pt]
Correction et amélioration & Traiter les anomalies et affiner l'ergonomie & Version corrigée et stabilisée \\
\addlinespace[2pt]
Finalisation & Consolider et documenter la solution & Application finalisée et documentation \\
\bottomrule
\end{tabular}
\end{center}

\section{Conclusion}

Ce chapitre a permis de présenter le cadre général du stage, l'organisme d'accueil ainsi que le contexte dans lequel s'inscrit le projet. Le fonctionnement de l'association, son organisation et les principaux besoins liés à la gestion des abonnés et des projets ont également été présentés.

Le sujet du stage, sa problématique ainsi que l'organisation du travail ont ensuite été définis. Le projet étant réalisé individuellement et à distance, une démarche progressive a été adoptée, allant de l'analyse des besoins jusqu'au développement, aux tests et à la finalisation de la solution.

Le chapitre suivant sera consacré à l'étude préalable du projet. Il permettra d'approfondir l'analyse des besoins, d'identifier les différents acteurs et de définir les fonctionnalités attendues de la solution.